% Compile with: uplatex file.tex && dvipdfmx file.dvi
%           or: latexmk -uplatex file.tex
\documentclass[dvipdfmx]{article}
\usepackage{amsmath, amssymb}
\usepackage{tikz}

\title{scenario\_free\_electron\_fermi\_surface}
\date{\today}

\begin{document}
\maketitle

Consider the free-electron-model dispersion. Prepare the wavevector $k$, electron mass $m$, chemical potential $\mu$, and reduced Planck constant $\hbar$ as scalar parameters.

Define k, m, $\mu$, $\hbar$ as scalars

\[
  k, m, \mu, \hbar
\]

Define the free-electron energy dispersion $\xi_k=\hbar^2k^2/(2m)-\mu$ (energy measured relative to the chemical potential).

$\xi_k = \hbar^2 k^2 / (2 m) - \mu$

\begin{equation}
  \xi_k = \hbar^2 k^2 / (2 m) - \mu
  \label{eq:E1}
\end{equation}

Simplify the right-hand side into a tidier form (no physical change).

Simplify the right-hand side of E1

\begin{equation}
  \xi_k = \frac{ - 2 m \mu + k^{2} \hbar^{2}}{2 m}
  \label{eq:E2}
\end{equation}

Fix the Fermi-surface condition $\xi_k=0$ (the energy coincides with the Fermi level) as an equation, by setting the right-hand side to 0.

Set the right-hand side of E1 to 0

\begin{equation}
  \hbar^2 k^2 / (2 m) - \mu = 0
  \label{eq:E3}
\end{equation}

Solve this equation for the wavevector $k$ to obtain the Fermi wavevector $k_F$.

Solve E3 for k

\begin{equation}
  k = \sqrt{2} \sqrt{m \mu}/\hbar
  \label{eq:E4}
\end{equation}

Substitute the Fermi wavevector just found back into the original expression for $\xi_k$, to confirm that the energy is identically zero on the Fermi surface.

Substitute E4 into E1

\begin{equation}
  \xi_k = \left( - \mu + \frac{\left( \frac{\sqrt{2 m \mu}}{\hbar} \right)^{2} \hbar^{2}}{2 m}\right)
  \label{eq:E5}
\end{equation}

Simplify the substitution result to make the zero explicit.

Simplify the right-hand side of E5

\begin{equation}
  \xi_k = 0
  \label{eq:E6}
\end{equation}

To find the effective mass $m^*$, take the second derivative of the dispersion $\xi_k$ with respect to the wavevector $k$.

Differentiate E1 twice with respect to k

\begin{equation}
  \frac{d^{2}}{dk^{2}}\left(\xi_k\right) = \frac{\hbar^{2}}{m}
  \label{eq:E7}
\end{equation}

From the second derivative just obtained, define the effective mass $m^*=\hbar^2/(d^2\xi_k/dk^2)$ (for the free-electron model this should equal $m$, the bare mass).

$m_star = \hbar^2 / E7$

\begin{equation}
  m_star = m
  \label{eq:E8}
\end{equation}

\end{document}